\section{Background}
\label{sec:background}

\subsection{The DocVQA 2026 benchmark}

Document Visual Question Answering asks a system to answer natural-language questions about a document presented as page images rather than as parsed text \citep{docvqa}. The ICDAR 2026 DocVQA benchmark targets the hard end of this task. Its documents are long, multi-page, and high-resolution, and they span eight content categories --- including dense business reports, scientific posters, maps, and engineering drawings --- so that answering a single question often requires locating the relevant page and region before reading it. The benchmark provides no training split: the validation set holds 25 documents and 80 questions, and the test set holds 48 documents and 160 questions. Models must therefore be developed without in-distribution supervised data, which places weight on transfer from related corpora and on the design of the inference-time agent.

Answers are scored with the Average Normalized Levenshtein Similarity (ANLS), introduced for scene-text VQA to give partial credit for near-miss strings while penalizing substantive errors \citep{stvqa}. ANLS is one minus the normalized Levenshtein edit distance between a prediction and a reference, averaged over questions. The 2026 protocol applies it in a binary form: a per-question similarity at or above a threshold of 0.9 counts as a correct answer and anything below it as wrong, so the reported score is effectively a thresholded accuracy. This binary ANLS@0.9 governs every number in this report; the evaluation protocol is given in full in \S\ref{sec:setup}.

Two properties of the documents shape the difficulty. First, much of the answer-bearing content is visual rather than textual --- figures, plots, schematics, and map labels --- where optical character recognition is unreliable or simply inapplicable, so a text-only pipeline cannot reach the relevant evidence. Second, many questions are multi-hop or arithmetic, requiring evidence to be gathered from several locations and then combined or computed rather than read off a single span. Both properties favor a system that can perceive document regions selectively and reason over what it perceives, rather than one that reads a fixed transcription of the page.

\subsection{Code-REPL agents for document VQA}

The agents studied here pair a frozen vision--language model (VLM) for perception with a code-capable language model (LM) as the reasoner. The reasoner does not view the document pixels directly; instead it writes code that issues perception requests against the VLM --- cropping, zooming, and querying specific page regions --- receives the returned observations into its context, and continues reasoning until it commits a final answer. Perception is thus \emph{active}: rather than consuming a single fixed rendering of the document, the reasoner decides what to look at next based on what it has already seen, iteratively, and on demand. This active visual perception is what lets the agent reach content OCR cannot and chain evidence across pages. The VLM stays a frozen external endpoint throughout; what varies between agents is the reasoning LM and the harness that connects it to perception.

These agents form a family. A harness exposes the same perception interface --- most importantly a code REPL in which the reasoner programs its visual queries --- but members differ in how they manage the growing context and which tools they surface; we refer to the members built around a code REPL and \texttt{batch\_look}-style perception as the REPL active-perception scaffold family. A prior competition entry, submitted under the name Perceive--Reason--Code, instantiates one member: a 27B-parameter reasoner driving such a harness over a frozen VLM, placed in the 8B--35B parameter tier of the leaderboard. That entry is one instance of the family, not a self-evidently best design, and it marks the prior-work boundary for our study.

What is reused from that prior work is the scaffold --- the active-perception harness and its tool surface. What is new here is twofold. First, a controlled \emph{harness evaluation}: rather than taking the active-perception design as given, we measure it against alternative agent designs to isolate which structural choices drive accuracy (\S\ref{sec:harness}). Second, \emph{weight-level training} of a small ($\leq$8B) reasoning LM inside the trainable member of the family, asking whether fine-tuning can lift a small agent further rather than deploying a larger frozen reasoner zero-shot (\S\ref{sec:training}). The scaffold is inherited; both the evaluation of it and the training within it are our contributions.
